## 5. Results and Evaluation — planning notes only, not a draft

Structure mirrors the methodology chapter (predictive / Avenue 1 / Avenue 2). Nothing here is a finished
section — it's a checklist of what each part needs, tagged `[exists]` (already computable from current
code/outputs, just needs the plot/table drawn) vs `[missing]` (needs new analysis code) vs `[partial]`.
Methodology describes the plan; this describes what actually happened when the plan was run — draft
properly later, this is a placeholder so nothing gets forgotten.

### 5.1 Predictive models (baselines, DNN, PINN, env-conditioned variants)

**Accuracy — does it answer the question (Q1/Q2)**
- `[exists]` MAE/RMSE/R² per model, per split, per cohort — scattered across per-run CSVs and `run_logs/`; needs one consolidated comparison table/plot, doesn't exist as a single artifact yet.
- `[exists]` Physics-weight ablation curve (accuracy vs. `physics_weight`/`trajectory_weight`) — numbers are in the experiment log; not yet plotted.
- `[missing]` A single "does physics help" summary figure (e.g. paired bars: DNN vs PINN vs zero-physics-control, per split) — the numbers exist, the figure doesn't.

**Capacity / training / overfitting**
- `[partial]` Per-epoch train/val loss is logged per run; a robustness check (loosened early stopping, ~2x training) exists for the no-env comparison and found DNN/4survey has a persistent overfitting climb that only softens with more training. Not yet compiled into a cross-model overfitting comparison plot (train-val gap vs. epoch, one line per model).
- `[missing]` Learning-curve analysis (accuracy vs. training-set size) — not found anywhere in the codebase; would need new runs, so this is training, not just re-evaluation — check before assuming it's "quick."

**Variance / error**
- `[partial]` Seed variation exists for the no-env physics-weight sweep (multiple seeds per setting) — not yet turned into a box/violin plot of R² across seeds per model.
- `[missing]` No seed-variation runs found for the env-conditioned models (E6/E7/E9) beyond the 5 CV folds themselves.

**Residual analysis — did the model work, outliers, does the next stage still see the same signal**
- `[exists]` Global Moran's I + LISA on CR-residual and on DNN/PINN/env-terrain residuals, cross-model comparison table already computed (`av1_spatial_autocorrelation_terrain.ipynb`) — CR residual has the *highest* spatial autocorrelation of all 6 models (I≈0.1223, range≈4,284m). This is the direct answer to "does the unexplained residual carry forward" for the predictive stage → Avenue 1's starting point. Needs the actual map/plot drawn (residuals plotted spatially, coloured by LISA cluster type), not just the summary numbers.
- `[missing]` No outlier-flagging analysis found for the predictive models specifically (as opposed to Avenue 2's disturbance-based exclusion, which is a different pipeline). Worth deciding: define "outlier" as top-N residual magnitude, or a formal threshold (e.g. >2-3 std)?
- `[missing]` Residual-vs-covariate plots (residual vs. age, vs. height, vs. survey interval, vs. management group) — not found; standard diagnostic, currently absent.

**Architecture diagrams** — described in methodology §4.2 (no-env DNN/PINN, `dnn_env_terrain`, `pinn_env_terrain`/`pinn_env_terrain_k`); actual figures still to be drawn from that text.

### 5.2 Avenue 1 (pooled CR residual attribution, `mean_cr_residual`)

- `[exists]` NLME variance-explained (compartment random intercept) — a real, if small, quantification of "how much of the residual is systematically explained" once terrain/wind fixed effects are added.
- `[exists]` Moran's I/LISA on `mean_cr_residual` itself (see 5.1) — this IS the "did the model work" residual check for the pooled CR curve.
- `[exists]` Paired fold-level significance test for `terrain_and_wind_only` vs `all_environmental` (computed 2026-08-09, in experiment log + `outputs/spatial_block_kfold/cr_residual_environmental_paired_significance_*.csv`).
- `[missing]` Outlier plot/list: which plots or compartments have the most extreme `mean_cr_residual`? Not yet produced.
- `[missing]` "Residual of the residual" check: after Elastic Net/XGBoost/NLME fit `mean_cr_residual`, does *their* residual still show spatial structure? Not found in the code search — this is the direct Avenue-1 analogue of the Avenue-2 check in 5.3, and is currently a real gap, not just an undrafted plot.
- `[missing]` Overfitting/variance for the attribution models themselves (EN/XGBoost train vs. test gap per fold, XGBoost early-stopping behaviour) — fold-level R² exists; a training-curve or train-vs-test-gap view does not.

### 5.3 Avenue 2 (per-plot growth-curve deviation, `local_y_max_difference`)

**Does it answer the question / accuracy**
- `[exists]` Pooled five-fold R² for EN/XGBoost/GNNWR across all four static scopes (log + CSVs).
- `[exists]` Paired significance test (Wilcoxon/t-test) and cluster-bootstrap CI, GNNWR vs. EN/XGBoost — already found the lead is not statistically confirmed (§4.8/§4.9 of the methodology chapter cite this).

**Capacity / training / overfitting**
- `[missing]` No training-curve or overfitting diagnostic found specifically for GNNWR (it's a gradient-trained neural model, so this applies) — worth checking whether `gnnwr_check.py` logs per-epoch loss at all.
- `[exists]` XGBoost/EN don't have a meaningful "overfitting" curve in the neural-network sense, but train-vs-test R² gap per fold is a direct analogue — computable from existing per-fold outputs, not yet compiled into a table.

**Variance / error**
- `[exists]` Per-fold R² spread already reported (e.g. terrain_wind fold R² range 0.011-0.253) — flagged in the log as "not a footnote," i.e. already recognised as important; needs a box-plot rendering.
- `[missing]` Repeated-seed CV (more fold-assignment seeds) for tighter significance — identified as the fix for the n=5 power problem (methodology §4.9), not run; would need new compute, so flag as future work, not a quick re-evaluation.

**Residual analysis — did the model work, outliers, carry-forward**
- `[exists]` Disturbance-based exclusion (`clearfell_like`, `measurement_inconsistent`) — this *is* the primary outlier-handling step for Avenue 2, done before fitting, not after.
- `[exists]` Residual spatial autocorrelation across all 5 models including GNNWR (KNN k=8, Moran's I) — found residual Moran's I stays high (~0.66-0.74) even for the best model (GNNWR) — the direct, already-computed answer to "does unexplained structure carry forward," stronger evidence than Avenue 1 currently has for the same question. Needs the spatial map drawn, not just the summary statistic.
- `[exists]` SHAP (`explain_signal.py`) and gain-based feature-importance rank stability across folds — both already computed, need figures (SHAP summary plot, rank-stability table/plot).
- `[missing]` Outlier plot for `local_y_max_difference` itself (beyond the disturbance-flag exclusion) — e.g. which retained, non-excluded plots still have extreme target values, and do they cluster spatially or by species/yield class?

### 5.4 Cross-cutting, not yet assigned to a subsection

- `[missing]` A single figure/table making the "three separate claims" framework (methodology §4.9) visually concrete: implementation-ran-correctly vs. behaviour-changed vs. prediction-improved, side by side per model — currently only stated as prose.
- `[missing]` Compute/runtime comparison across models — mentioned as "report only if measured consistently" in the evaluation rules; not confirmed whether wall-clock time was logged consistently enough to use.
